\scanpage{17}
Si par exemple on prend la $W$ définie par (*), \deleted{$W(f)<\infty\ \unsure{\Longleftrightarrow}\ f\in L^p(M)$, \unsure{et}} on a $W(f)=\int_M|f|^p$, donc $W(f)<+\infty$ équivaut à $N_p(f)<+\infty$.

3. \uline{Inégalités liées aux \deleted{fonctions de Polya} \inserted{fonctions de Weyl.}}

\uline{Prop 4} \inserted{Soit $0<\omega<\infty$.} Soient $f,g$ deux fonctions \inserted{sommables} décroissantes sur $]0,\omega[$ (\unsure{essentiellement finies|essentiellement bornées}), telles que
\[
(1)\qquad \boxed{\int_0^t f\leq\int_0^t g}\quad\inserted{\text{pour }0\leq t\leq\omega.}
\]
Alors pour toute fonction de Polya \inserted{croissante} $P$ sur $L^\infty(]0,\omega[)$, on a
\[
(2)\qquad P(f)\leq P(g)
\]

En vertu de N° 2, (1), il suffit de prouver que
\[
\int_0^\omega f\varphi\leq\int_0^\omega g\varphi
\]
pour toute $\varphi$ \uline{positive décroissante} sommable sur $(0,\omega)$.

Or cela résulte facilement de (1) par intégration par parties. En posant $F(t)=\int_0^t f$, on a
\[
(3)\qquad\int_0^\omega f\varphi=[\varphi F]_0^\omega+\int_0^\omega F(-d\varphi)=\varphi(\omega)F(\omega)+\int_0^\omega F(-d\varphi)
\]

\uline{Remarque} \inserted{(1) La réciproque est d'ailleurs triviale, si on remarque que pour tout $t$, $f\mapsto\int_0^t\varphi_f$ est une fonction de Polya sur $L^\infty(]0,\omega[)$.}
On voit sur (3) que si $F(\omega)=G(\omega)$, i.e. $\int_0^\omega f=\int_0^\omega g$, il n'est pas nécessaire de supposer $\varphi$ positive dans la conclusion de la prop 4. Nous ne nous \unsure{servirons} pas de cette remarque, qui donne des
\mspage{18}
